\section{Physics-Consistent Generative Inverse Design}

\subsection{Simulation remains the source of physical authority}

The starting point is a simulation-grounded dataset of freeform binary structures paired with their angle-resolved transmission responses. This choice follows directly from the physics of the problem. For Fourier-operator metasurfaces, the quantity that ultimately matters is the shape of the response field, and that field must be judged by Maxwell-consistent electromagnetic evaluation. In the present project, RCWA is used to annotate each structure over a wavelength--angle grid for a fixed physical stack and material system. The computational burden is substantial, but reducing the target to a weaker scalar proxy would erase the very structure that distinguishes a genuine operator from a locally adequate device \cite{moharam1981,moharam1995}.

The computational burden of this supervision is precisely what makes freeform inverse design difficult. Each candidate must, in principle, be evaluated over a finite angular and spectral aperture, and the resulting response must be judged with operator-aware criteria rather than with a single transmission number. The design loop adopted here therefore does not attempt to replace physics with a learned model. Simulation anchors both ends of the procedure: it provides the supervision signal during training, and it supplies the final decision rule after candidate generation.

\subsection{Forward surrogates address screening burden, not physical closure}

Dense RCWA evaluation for every generated candidate is expensive enough that some learned approximation becomes practically useful. The forward surrogate serves that role. It is trained to approximate the mapping from structure to response field and is used to reduce screening burden by concentrating attention on candidates that are more likely to satisfy the target. The surrogate is therefore best understood as a solver-allocation tool: it tells us where full-wave effort is worth spending, but it does not decide what counts as a valid device.

This distinction matters because surrogate fidelity and physical fidelity are not interchangeable. A learned predictor may recover broad trends in the response landscape while missing the local structure that determines whether target curvature, edge roll-off, or polarization leakage is acceptable. For operator metasurfaces, such differences are often the difference between real operator behavior and a visually plausible approximation near one nominal condition. The present code base reflects that principle directly: the forward model is trained with a physical-space reconstruction loss and used to guide candidate concentration, whereas final acceptance is deferred to RCWA evaluation. The surrogate reduces electromagnetic cost, but only partially and intentionally.

\subsection{Why the inverse stage must be generative}

The inverse map from response field to geometry is intrinsically multimodal. Even for a highly structured target, there is rarely one binary layout that deserves to be called the only correct solution. This ambiguity becomes more pronounced as the geometry class becomes freer and the target description becomes richer. In that setting, deterministic regressors tend to return an average over several latent design families, often producing weaker best-case behavior after physical verification. Generative models are attractive here for a narrower reason: they preserve the possibility that several distinct structures may satisfy the same optical constraint.

Diffusion models are well suited to this role because they support conditional sampling over complex, high-dimensional distributions and have already shown promise in photonic inverse design \cite{liu2018,so2019,ma2021,tanriover2023,zhang2023,zhang2024,hen2026}. Here the conditional input is the full target response field rather than a reduced descriptor. The model is therefore asked to infer geometry from a continuous optical specification defined over angle, wavelength, and polarization. That is where the response-space formulation and the generative view meet most naturally: the richer the target, the less plausible it becomes that one deterministic inverse image should suffice.

\subsection{Physics consistency is enforced by closed-loop verification}

The final step is post-generation physical verification and reranking. Generated candidates are re-evaluated by the electromagnetic solver, and task-specific scores are computed from the verified response rather than from the surrogate prediction alone. This is not a cosmetic correction added after a machine-learning model has already ``solved'' the design problem. It is the step that turns candidate generation into a physically trustworthy inverse-design procedure.

The need for reranking is evident in the project outputs themselves. Candidate sets often contain structures with similar nominal scores but meaningfully different behavior in edge regions, angular shoulders, spectral drift, or passive polarization leakage. The ranking scripts therefore separate shape agreement, center suppression, edge behavior, bandwidth preservation, and channel leakage rather than collapsing everything into a single scalar. A physics-consistent design loop must preserve candidate diversity long enough for this verified comparison to matter. That is why the present approach is best understood as \emph{physics-consistent generative inverse design}: candidate generation explores the feasible set, the surrogate reduces screening burden, and the electromagnetic solver closes the loop.
